% (c) Nikita Lisitsa, lisyarus@gmail.com, 2026

\documentclass[10pt]{beamer}

\usepackage[T2A]{fontenc}
\usepackage[russian]{babel}
\usepackage{minted}

\usepackage[most]{tcolorbox}
\tcbuselibrary{minted}

\usepackage{graphicx}
\graphicspath{ {./images/} }

\usetheme{metropolis}

\usepackage{adjustbox}

\usepackage{color}
\usepackage{soul}

\usepackage{hyperref}

\definecolor{codebg}{RGB}{29,35,49}

\setminted{fontsize=\footnotesize}

\setlength\partopsep{-\topsep}

\definecolor{blue}{rgb}{0,0,1}
\definecolor{red}{rgb}{1,0,0}

\usemintedstyle{tango}

\makeatletter
\newcommand{\slideimage}[1]{
  \begin{figure}
    \begin{adjustbox}{width=\textwidth, totalheight=\textheight-2\baselineskip-2\baselineskip,keepaspectratio}
      \includegraphics{#1}
    \end{adjustbox}
  \end{figure}
}
\makeatother

\title{Компьютерная графика}
\subtitle{Практика 6: render to texture}
\date{2026}

\setbeamertemplate{footline}[frame number]

\begin{document}

\frame{\titlepage}

\begin{frame}[fragile]
\frametitle{Практика 6}
\begin{itemize}
\item Модель можно двигать и вращать стрелками
\end{itemize}
\end{frame}

\begin{frame}[fragile]
\frametitle{Практика 6}
\slideimage{0.png}
\end{frame}

\begin{frame}[fragile]
\frametitle{Задание 1}
Создаём sampler, offscreen текстуру, view, и bind группу
\begin{itemize}
\item Параметры sampler'а: \mintinline{cpp}|address_mode = ClampToEdge|, \mintinline{cpp}|*_filter = Linear|
\item Формат текстуры -- \mintinline{cpp}|RGBA8Unorm|, размер -- \(\frac{width}{2}\times\frac{height}{2}\)
\item Bind группа использует уже имеющийся в коде \mintinline{cpp}|bindGroupLayout|, в ней 2 entry: sampler и texture view
\item Текстуру, view и bind группу нужно пересоздавать при изменении размеров экрана (так же, как буфер глубины)
\item \textbf{\alert{NB}}: Пока ничего не изменится, мы никак не используем эту текстуру :)
\end{itemize}
\end{frame}

\begin{frame}[fragile]
\frametitle{Задание 2}
Рисуем созданную текстуру
\begin{itemize}
\item \textbf{После первого render pass'а} создаём второй, тоже рисующий на экран, но \textbf{без буфера глубины}; \mintinline{cpp}|loadOp = Load|, \mintinline{cpp}|storeOp = Store|
\item Выставляем созданную bind группу под индексом 0: \mintinline{cpp}|render_pass.set_bind_group(0, bind_group)|
\item В коде уже есть \mintinline{cpp}|texturePipeline| -- выставляем ей \mintinline{cpp}|TextureImmediates| и рисуем \mintinline{cpp}|render_pass.draw(vertex_count = 6)|, чтобы в левой нижней четверти экрана появился прямоугольник
\item \textbf{\alert{NB}}: Текстура будет чёрной -- мы пока ничего в неё не рисовали
\end{itemize}
\end{frame}

\begin{frame}[fragile]
\frametitle{Задание 2}
\slideimage{2.png}
\end{frame}

\begin{frame}[fragile]
\frametitle{Задание 3}
Рисуем в нашу текстуру
\begin{itemize}
\item Изменяем настройку \textbf{первого render pass'а}, чтобы он рисовал в нашу текстуру
\item Придётся изменить формат пикселя при создании \mintinline{cpp}|mainRenderPipeline|
\item Придётся изменить размер буфера глубины (до совпадающего с размером текстуры)
\end{itemize}
\end{frame}

\begin{frame}[fragile]
\frametitle{Задание 3}
\slideimage{3.png}
\end{frame}

\begin{frame}[fragile]
\frametitle{Задание 4}
Создаём array-текстуру
\begin{itemize}
\item При создании нашей текстуры указываем размер по третьей координате равный \mintinline{cpp}|Z = 4|
\item Меняем одну texture view и одну bind группу на массивы из 4х текстур и bind групп, соответственно
\item Texture view с номером \(k\) будет ссылаться на \mintinline{cpp}|k|-ый слой array текстуры, для этого при её создании нужно указать \mintinline{cpp}|base_array_layer = k|
\item \mintinline{cpp}|k|-ая bind группа ссылается на \mintinline{cpp}|k|-ый texture view
\item Для рендеринга пока используем \mintinline{cpp}|0|-ую bind группу и \mintinline{cpp}|0|-ой texture view
\item \textbf{\alert{NB}}: Пока ничего не изменится, мы никак не используем оставшиеся 3 слоя
\end{itemize}
\end{frame}

\begin{frame}[fragile]
\frametitle{Задание 5}
Рисуем модель 4 раза в разные слои текстуры
\begin{itemize}
\item Заносим \textbf{первый render pass} в цикл: \mintinline{cpp}|k|-ая итерация рисует в \mintinline{cpp}|k|-ый слой текстуры (т.е. в \mintinline{cpp}|k|-ый texture view)
\item У каждой итерации делаем свой цвет очистки и свою матрицу проекции (одна перспективная, три ортографических с разных сторон)
\item Заносим рисование прямоугольника с текстурой на экран в цикл: \mintinline{cpp}|k|-ая итерация рисует \mintinline{cpp}|k|-ый слой (т.е. использует \mintinline{cpp}|k|-ую bind группу)
\item Прямоугольники должны вместе покрыть экран
\item \textbf{\alert{NB}}: В сумме должно получиться \textbf{5 render pass'ов}
\end{itemize}
\end{frame}

\begin{frame}[fragile]
\frametitle{Задание 5}
\slideimage{5.png}
\end{frame}

\begin{frame}[fragile]
\frametitle{Задание 6}
Добавляем color grading
\begin{itemize}
\item Добавляем в \mintinline{cpp}|TextureImmediates| поле \mintinline{cpp}|mode: i32| и устанавливаем его в номер текущей итерации
\item В шейдере \mintinline{cpp}|shader_texture.wgsl| при \mintinline{cpp}|mode == 1| делаем color grading: преобразуем выходной цвет как
\usemintedstyle{lightbulb}
\begin{tcblisting}{listing engine=minted, minted language=wgsl, minted options={fontsize=\scriptsize}, listing only, colback=codebg, colframe=codebg, rounded corners}
color = floor(color * 4.0) / 3.0;
\end{tcblisting}
\usemintedstyle{tango}
\item В противном случае возвращаем цвет без изменений
\end{itemize}
\end{frame}

\begin{frame}[fragile]
\frametitle{Задание 6}
\slideimage{6.png}
\end{frame}

\begin{frame}[fragile]
\frametitle{Задание 7}
Добавляем domain warping
\begin{itemize}
\item Добавляем в \mintinline{cpp}|TextureImmediates| поле \mintinline{cpp}|time: f32| и передаём туда время (оно уже есть в заготовке)
\item В шейдере \mintinline{cpp}|shader_texture.wgsl| при \mintinline{cpp}|mode == 2| делаем domain warping: перед чтением из текстуры искажаем текстурные координаты чем-нибудь зависящим от времени, например
\usemintedstyle{lightbulb}
\begin{tcblisting}{listing engine=minted, minted language=wgsl, minted options={fontsize=\scriptsize}, listing only, colback=codebg, colframe=codebg, rounded corners}
vec2(sin(in.texcoord.y * 50.0 + immediates.time) * 0.01, 0.0)
\end{tcblisting}
\usemintedstyle{tango}
\end{itemize}
\end{frame}

\begin{frame}[fragile]
\frametitle{Задание 7}
\slideimage{7.png}
\end{frame}

\begin{frame}[fragile]
\frametitle{Задание 8*}
Добавляем гауссово размытие
\begin{itemize}
\item В шейдере \mintinline{cpp}|shader_texture.wgsl| при \mintinline{cpp}|mode == 3| делаем гауссово размытие, например, с \mintinline{cpp}|N = 10| и \mintinline{cpp}|sigma = 6|
\item Если будет тормозить, можно взять значения поменьше
\item Код шейдера можно взять из слайдов лекции
\end{itemize}
\end{frame}

\begin{frame}[fragile]
\frametitle{Задание 8*}
\slideimage{8.png}
\end{frame}

\end{document}
